\subsection{Five even multipair packets with the actual two-step root exception}\label{proof:Q5-C182}
% Common local aliases for the former repeated C168.1--C168.3 prefaces.
Use the original-measure calculus of
Section~\ref{sec:odd-common-calculus} and the complete root-pair
certificates of Section~\ref{sec:odd-common:fourth}.
For a root \((1,2,v)\), write \(m=v-2\), \(G_m=E+zD\),
\(\beta=\beta_{v-1}=\int D\,d\nu\),
\(\beta_v=\int E\,d\nu\),
\(P_+=h_2f_v\), \(P_-=\bar f_2W_v\),
\(\ell_+=\ell_a\), and \(\ell_-=\ell_s\).
The two vectors, their signed means and full constants are
\eqref{eq:odd-common:root-standard}--\eqref{eq:odd-common:root-losses}.
At \(v=3\), \((E,D)=(1,0)\) and \(\beta=0\).
All subsequent direction choices and packet estimates retain their
stated integer and parameter ranges.
Here $n=v-2$. The complete finite unequal-pair mask proofs and all five terminal packets are in R29.3.2, R29.6.1--R29.6.2 of Section~\ref{proof:R2-C029}. The following payments retain the genuine positive root $P_+$ precisely at $L=2$; the negative root $P_-$ is used at all other lengths.
\paragraph{C182.0 Original parameters and the genuine two-step exception}

All original arms, equal-pair multiplicities, subdivisions, independent positive core lengths, edge counts and arbitrary label laws remain.

The internal root residual is (1, 2, v), v at least three. The terminal residual is
\[
 1\le a_1<b_1<a_2<b_2<\cdots<a_j<b_j,\qquad j\ge2.       \tag{C182.1}
\]
The positive original inner and outer lengths \(L,\allowbreak \allowbreak  M\) are independent, and arbitrary genuine equal pairs may remain at both roots.

% Replace a pure model-definition block only; keep local root/terminal data.
Use the original finite-space model and arm messages of
Section~\ref{sec:odd-common-calculus}.
When the original \(K\)-notation is used, \(K=-B\) and
\(g=K\one=-u\); in either notation
\(\delta_j=\|B^j\one\|_2^2=\|K^j\one\|_2^2\).
All means, norms, adjoints and spectral integrals use the same
original probability measure.
Fix the actual directions as
\[
 F=\prod_{i=1}^j h_{a_i}f_{b_i},\qquad
 P=\begin{cases}h_2f_v,&L=2,\\\bar f_2W_v,&L\ne2,\end{cases}
 \qquad H=pP T^LF.                                    \tag{C182.2}
\]
The root unit arm points outward and gives \(p\). Inner edges point terminal to root, outer edges root to receiver. If \(a_1=1\), its \(h_1=q\) is an actual terminal unit arm. The two root pairs correspond to \(L=2\) and \(L\ne2\), respectively.

Let
\[
 \lambda_j=\frac{27}{5}+\frac{64}{5}(j-2).
\]
For every listed original parameter and genuine equal-pair decoration,
\[
 \boxed{EH^4\le e^{-\lambda_jd}.}                       \tag{C182.3}
\]
For \(a_1=1\), strengthen the basic rate \(27/5\) to \(13/2\). Apply the Q5-C164 all-length outlet to obtain
\[
\begin{gathered}
I_q(T^MH)^3\le e^{-[2\lambda_j-c(M)]d}\le e^{-6d/95},\\
\quad I_q(X)=E[q^{-1/3}|X|^{8/3}],
\end{gathered}\tag{C182.4}
\]
Here \(c(1)=204/19,\allowbreak \allowbreak  c(2)=10,\allowbreak \allowbreak  c(3)=0,\allowbreak \allowbreak  c(4)=7\), and \(c(M)=6\) for \(M\ge5\).

For two such sides and the compatible negative zero-step center on an original five-branch path, all four positive core lengths remain independent and
\[
 \boxed{Z_D\le e^{-3d/190}\le1.}                       \tag{C182.5}
\]
If both terminal minima are one, the stronger exponent is \(43/76\).

\paragraph{C182.2 One actual exponent and every finite propagation}

Suppose \(F\le e^{R-\Gamma}\), and the same pointwise domain has \(e^R\le1+R+cR^2/2\). Put
\[
 \sigma=ER^2,\quad D_k=EB^kR,\quad
 \delta=D_0+c\sigma,\quad\epsilon=\delta+D_1.
\]
For \(|R|\le t<3\), the positive Taylor bound \(n!\ge2\cdot3^{n-2}\) permits \(c=(1-t/3)^{-1}\); negative \(R\) has integral coefficient at most one. Thus the bounds \(t=3/5,\allowbreak \allowbreak  1/3,\allowbreak \allowbreak  1,\allowbreak \allowbreak  3/4\) allow \(c=5/4,\allowbreak \allowbreak  6/5,\allowbreak \allowbreak  3/2,\allowbreak \allowbreak  4/3\). Positivity and \(T(R^2)\le2\sigma\) give
\[
\begin{gathered}
T^LF\le e^{-\Gamma}Y_L,\quad\\
Y_L=\bar f_L+\delta \bar f_{L-1}+B^LR+\\
\sum_{i=1}^{L-1}D_i \bar f_{L-1-i}>0.
\end{gathered}\tag{C182.14}
\]
This is finite Duhamel induction. The scalar \(\delta\) may be negative, but \(Y_L\) dominates the actual positive message. Combine \(\log Y_L\le Y_L-1\) with the complete root envelope and the original \(p\) curvature. The curvature coefficient \(c_0\) is distinct from the Taylor coefficient.

At \(L=1\), use \(V=R-BG_n(B)u\), \(c_0=\beta\), giving
\[
\begin{gathered}
\log EH^4\le-4\Gamma+4\epsilon-4\ell_-\\
+(17/4)\|V\|^2+(10/3)\beta^2.
\end{gathered}\tag{C182.15}\label{eq:five-even-multipair:L1}
\]
At \(L=2\), retain the actual positive-root exception. Its vector is \(V=[1-BG_n(-B)]u+BR\), and its full curvature coefficient is \(c_0=\beta+\delta\):
\[
\begin{gathered}
\log EH^4\le-4\Gamma+4\epsilon-4d+4D_2-4\ell_+\\
+(17/4)\|V\|^2+(10/3)(\beta+\delta)^2.
\end{gathered}\tag{C182.16}\label{eq:five-even-multipair:L2}
\]
At \(L=3\), use the negative root, \(V=[G_2-BG_n]u+\delta u+B^2R\), with full \(c_0=\beta+\epsilon\):
\[
\begin{gathered}
\log EH^4\le-4\Gamma+4\epsilon+4D_2+4D_3-4\delta d\\
+J_3+(10/3)(\beta+\epsilon)^2,
\end{gathered}\tag{C182.17}\label{eq:five-even-multipair:L3}
\]
Here \(J_3=-8d-4\ell_-+17\|V\|^2/4\), retaining this same actual \(\delta,\allowbreak \allowbreak  R\).

For \(L\ge4\), put \(\gamma_k=\beta_k-D_{k-1}\). If every \(k\ge3\) has \(\gamma_k\ge hd\), finite convolution retains the original final \(\bar f_0,\allowbreak \allowbreak  \bar f_1\) terms:
\[
 Y_L\le S_L+B^LR+2\delta_++2(D_1)_+-hd(2+u).
\]
The actual mean also contains \(-\beta_{L+1}+D_L=-\gamma_{L+1}\le-hd\). Set \(V=[G_{L-1}-BG_n]u+B^{L-1}R\), \(c_0=\beta-hd\), obtaining
\[
\begin{gathered}
\log EH^4\le-4\Gamma+8\delta_++8(D_1)_+-12hd-4\ell_-\\
+(17/4)\|V\|^2+(10/3)(\beta-hd)^2.
\end{gathered}\tag{C182.18}\label{eq:five-even-multipair:Lge4}
\]
The gamma proof must include every finite index, especially \(L+1\). Every discarded path message has a proved nonnegative subtracted coefficient.

\paragraph{C182.2.1 Three finite common bases}

For every n at least one, the two-step positive-root base satisfies
\[
 |1-zG_n(-z)|^2=(1-xD)^2+xE^2\le1+x<(19/16)^2.      \tag{C182.19}
\]
The negative-root three-step base \(Q=G_2-zG_n\) has modulus below \(23/16\), and the stronger joint payment
\[
 -8d-4\ell_-+(17/4)\|Q(B)u\|^2\le-27d/4.         \tag{C182.20}
\]
At \(n=1\), the left side equals \(-31d/4\). For \(n\ge2\), the frequency expression decreases in both \(E,\allowbreak \allowbreak  D\): the \(D\)-derivative is at most \(-4+9x/2+17x^2/2<0\), the \(E\)-derivative at most \(-4\). Set both to \(1-x\), obtaining
\[
 -47/4+(25/2)x-x^2/4-(17/4)x^3+(17/4)x^4.
\]
This increases on \([0,\allowbreak \allowbreak  \theta]\), with endpoint \(-28866559/4194304<-27/4\), a margin of \(555007/4194304\).

For every m at least three and n at least one, the long base satisfies
\[
\begin{gathered}
|G_m-zG_n|^2\le(1+x^2)^2+x^3\\
\le1493553/1048576<36/25.
\end{gathered}\tag{C182.21}
\]
The real coordinate of \(G_m\) is at most \(1-x+x^2\), and the imaginary-coordinate difference has absolute value at most \(x\), including the actual \((1,\allowbreak \allowbreak  0)\) endpoint. Every finite length is covered.


\paragraph{C182.4 Packet E2: terminal minimum two}

Take the actual finite mask \(R=\sum_iR_{a_ib_i}=-B^2P(-B)u\), with \(P_0=1\), and \(\Gamma\ge607d/160+\Delta\), \(\Delta=16(j-2)d/5\). Its complete zero-one-mask, coordinate, signed-mean and all-index gamma proofs are parts R29.6.1--R29.6.2 of Section \ref{proof:R2-C029}. In particular,
\[
\begin{gathered}
\sigma\le(5408/16245)d<d/3,\\
\quad\|BR\|^2\le(2197/16245)d<(37/100)^2d.
\end{gathered}\tag{C182.27}
\]
The row bound gives \(|R|^2\le C_*\theta d<9/25\); use \(c=5/4\). For \(R=(\alpha+z\beta_0)u\), the actual coordinates satisfy
\[
 x/2\le \alpha\le6x/5,\qquad -6x/5\le beta0\le x/2.
\]
The forced original \(K^3\) term gives \(D_1\le0\); the complete same-sign moment tails give
\[
\begin{gathered}
-d/5\le\delta\le(4888/5415)d<(91/100)d,\\
\quad |\epsilon|\le7d/10,
\end{gathered}\tag{C182.28}
\]
\[
 D_2\le d/12,\quad D_3\le d/5,\quad D_k\le d/12\ (k\ge4).
\]
The positive delta coefficient is exactly \(416/855+(5/4)(5408/16245)=4888/5415\). Epsilon retains the same-model inequality \(D_0+D_1\le13\delta_2/19\), giving its positive bound; its lower bound uses \(D_0\ge-d/5,\allowbreak \allowbreak  D_1\ge-d/2\). All later positive moment tails start at \(\delta_4\), with same-sign indices separated by two. Hence \(\gamma_3\ge11d/12,\allowbreak \allowbreak  \gamma_4\ge4d/5\), and all later gammas at least \(49d/96\); take \(h=1/2\).

At \(L=1\), the coordinates of \(V\) are \(\alpha+xD,\allowbreak \allowbreak  \beta_0-E\), with square at most \((11x/5)^2+x(1+6x/5)^2\le16/9\), using \(x\le5/12\). The complete first-row payment, with \(\beta^2\le d^2\), gives basic loss
\[
 607/40-14/5+4-68/9-10/9=185/24.                    \tag{C182.29}
\]

At \(L=2\), the positive-root base and \(BR\) bound give \(\|V\|<623\sqrt d/400\). The actual coefficient obeys \(|\beta+\delta|\le191d/100\). With \(d\le1/3\), the loss is at least
\[
\begin{gathered}
\frac{385}{24}+\frac{969}{256}\\
-\frac{6598193}{640000}-\frac{36481}{9000}\\
>\frac{27}{5}.
\end{gathered}\tag{C182.30}
\]
The left side has numerator \(31470923\) over \(5760000\), while \(27/5\) has numerator \(31104000\); the margin is \(366923\). The actual delta constant is paid through its original curvature.

At \(L=3\), the jointly paid base has modulus bound \(23/16\). The perturbation \(\delta u+B^2R\) has norm coefficient at most \(91/300+1/4=83/150\), adding norm cost
\[
 (17/4)[(23/8)(83/150)+(83/150)^2]
 =2902427/360000<81/10.
\]
The additional mean cost is at most \(21d/5\), and \(|\beta+\epsilon|\le17d/10\) costs at most \(289d/90\) in curvature. Thus the basic loss is
\[
 607/40+27/4-81/10-21/5-289/90=2309/360>32/5.       \tag{C182.31}
\]

For \(L\ge4\), \(\|B^{L-1}R\|<3\sqrt d/20\), since \(\theta^3/3<9/400\). The full vector has norm below \(27\sqrt d/20\), \(|\beta-d/2|\le d/2\), and \(2\delta_++2(D_1)_+\le91d/50\). The basic loss is at least
\[
 607/40+68/25-12393/1600-5/18>49/5.                 \tag{C182.32}
\]
All four lengths pay \(27/5\), with an additional \(64/5\) for each later terminal pair.

\paragraph{C182.5 Packet E3: terminal minimum at least three}

Take the same finite interval mask, now with \(\Gamma\ge767d/160+\Delta\). Its full derivation is the E3 packet in R29.6.2, including every finite signed moment. The lowest original power is four, giving
\[
 \|R\|<(3/8)\sqrt d,\quad |R|<1/3,
\]
Use \(c=6/5\) on the full domain. The same-sign moment tails give
\[
 -d/5\le\delta\le2d/3,\quad |\epsilon|\le7d/8,\quad D_1\le d/5,
\]
\[
 D_2\le d/12,\quad D_3\le d/30,\quad D_k\le d/12\ (k\ge4). \tag{C182.33}
\]
Explicitly \(\sigma\le2197d/16245\), and \(\delta\le52702d/81225<2d/3\). Since \(D_1\ge-d/12\), \(\epsilon\ge-17d/60\), while its upper bound is \(13d/15<7d/8\). The gamma bounds are \(11d/12,\allowbreak \allowbreak  29d/30\) at the first two indices and \(49d/96\) thereafter, allowing \(h=1/2\).

At \(L=1\), \(\|BG_n(B)u\|<4\sqrt d/5\) gives \(\|V\|<47\sqrt d/40\), with basic loss \(787/40-37553/6400-10/9>12\). At \(L=2\), \(\|BR\|<\sqrt d/4\) and the positive-root base give \(\|V\|<23\sqrt d/16\), \(|\beta+\delta|\le5d/3\), and loss \(767/40-7/2+4-1/3+969/256-8993/1024-250/81>11\).

At \(L=3\), the perturbation coefficient is at most \(2/9+1/6=7/18\), adding norm cost \(27965d/5184<11d/2\). The actual additional mean costs at most \(127d/30\); \(|\beta+\epsilon|\le15d/8\) gives curvature cost at most \(125d/32\). Hence the basic loss is \(767/40+27/4-11/2-127/30-125/32>12\).

For \(L\ge4\), the tail norm is below \(\sqrt d/10\), so the full vector coefficient is below \(13/10\). The positive scalar terms satisfy \(2\delta_++2(D_1)_+\le26d/15\), and the curvature cost is at most \(5d/18\). The loss is \(767/40+46/15-2873/400-5/18>14\). Thus every finite parameter in E3 pays the common rate and pair-count increment.


\paragraph{C182.7 All unit and two-step connectors for minimum one}

The three packets I--III are completely derived, including their forced missing exponents, masks, scalar losses and signed moments, in R29.6.1--R29.6.2. They are I: \(a_1=1,\allowbreak \allowbreak  b_1\ge3\); II: \(a_1=1,\allowbreak \allowbreak  b_1=2,\allowbreak \allowbreak  a_2\ge4\); III: \(a_1=1,\allowbreak \allowbreak  b_1=2,\allowbreak \allowbreak  a_2=3\). Their basic gamma lower bounds are \(607d/160,\allowbreak \allowbreak  16d/5,\allowbreak \allowbreak  3d-2\delta_2\), respectively, plus \(\Delta\). Write \(R=(xV+\xi)+zU\), with \(\xi=0\) in I--II and the original global \(\xi=\tau\) in III. All use the negative root at one step and the actual positive root at two.

\paragraph{C182.7.1 Packet I}

At \(L=1\), the coordinates are \(x(V+D),\allowbreak \allowbreak  U-E\), with absolute values at most \(11x/5,\allowbreak \allowbreak  5/8\). The square is below \(121/144+325/2048=18413/18432<1\). The basic loss is \(607/40-3+4-17/4-10/9>10\).

At \(L=2\), the coordinates are \(1-xD-xU,\allowbreak \allowbreak -E+xV\). The first is positive and the second negative, since \(xV\le416/855<19/32\le E\). Both absolute values are at most \(1-x/2\), so their square is at most \((1+x)(1-x/2)^2\le1\). The actual curvature coefficient has absolute value at most \(7d/4\), costing \(245d/72\). The basic loss is \(607/40-3+4-2+969/256-17/4-245/72>10\).

\paragraph{C182.7.2 Packet II}

At \(L=1\), the real-coordinate absolute value is below \(1/2\), and imaginary coefficient below \(5/8\); the squared modulus is below \(5/12\). The basic loss is \(396/25-85/48-10/9>12\).

At \(L=2\), the real coordinate is positive and at most one, and the negative imaginary coefficient has absolute value at most \(6/5\). The square is at most \(1+36\theta/25=317/200<8/5\). The actual curvature coefficient has absolute value at most \(6d/5\), costing \(8d/5\). The loss is \(64/5-24/25+4-2+969/256-34/5-8/5>9\).

\paragraph{C182.7.3 Packet III}

At \(L=1\), the real coordinate \(\tau+x(V+D)\) has positive bound \(205/288<5/7\) and negative absolute bound \(C_\beta x^2<1/5\). The imaginary coefficient \(U-E\) has absolute value at most \(1/2\). Thus the square is below \(5/8\), giving loss
\[
 227/20-85/32-10/9>15/2.                            \tag{C182.41}
\]

At \(L=2\), the real coordinate lies in \([1-2x,\allowbreak \allowbreak  1-x/2]\); the imaginary coefficient \(-E+\tau+xV\) is negative, since \(\tau+xV\le2/9+1/12<19/32\), and has absolute value at most \(6/5\). Its square is less than \((1-x/2)^2+36x/25<5/4\), by \(1+143/800+169/4096<5/4\). With actual constant coefficient at most \(23d/20\), curvature costs \(529d/360\), so the basic loss is
\[
 35/4-7/5+4-13/8+969/256-85/16-529/360>13/2.        \tag{C182.42}
\]
The actual \(D_2\le \delta_2\le\theta d\) is used; the global \(\delta_2\) is constant.

\paragraph{C182.8 The three-step and every longer connector for minimum one}

\paragraph{C182.8.1 The full three-step continuous-frequency comparison}

The coordinates of the same full vector are
\[
 A_3=1+xD-x\xi-x^2V+\delta,\qquad C_3=1-E-xU,
\]
\[
 J_3=\int[-8-4E-4(1+x)D+(17/4)(A_3^2+xC_3^2)]d\nu. \tag{C182.43}
\]
The scalars \(\xi\), either zero or global \(\tau\), and \(\delta\) retain their actual global values. The real coordinate is positive in all three cases, so its proved delta upper bound may be used.

In I, use \(\delta\le1/4,\allowbreak \allowbreak  V\ge1/2\). The \(D\)-derivative is at most \(-4+53x/8+17x^2/2-17x^3/4\), increasing to \(-493/6912<0\) at \(5/12\). In II and III use \(\delta\le1/20,\allowbreak \allowbreak  V\ge-6x/5,\allowbreak \allowbreak  \xi\ge0\). The derivative is at most \(-4+197x/40+17x^2/2+51x^4/5\), with numerator \(-17085\) over \(103680\) at \(5/12\). The \(E\)-derivative is at most \(-4+51x^2/5<0\) in every packet. Thus set \(E=D=1-x\) for \(n\ge2\), retaining \((1,\allowbreak \allowbreak  0)\) separately at one.

In I, the later-root real coordinate is at most \(17/12\), imaginary coefficient at most \(1/5\), and square below \(49/24\), giving \(J_3\le-3511d/256+833d/96<-5d\). At one, the square is below \(5/3\), giving \(J_3\le-59d/12\); uniformly use \(-19d/4\).

In II, the later-root real bound is \(101/75\), imaginary bound \(1/12\), square below \(11/6\), giving \(-3511d/256+187d/24<-23d/4\). At one, the corresponding bounds are \(11/10,\allowbreak \allowbreak  1/2,\allowbreak \allowbreak  4/3\), with the same joint conclusion.

In III, the later-root bounds \(83/60,\allowbreak \allowbreak  1/5,\allowbreak \allowbreak  39/20\) give \(J_3\le-3511d/256+663d/80=-6947d/1280<-27d/5\). At one, bounds \(17/15,\allowbreak \allowbreak  x,\allowbreak \allowbreak  11/8\) give \(-197d/32<-27d/5\). These coordinate comparisons use \(x-x^2\le1/4,\allowbreak \allowbreak  C_\beta x^3<1/12\), and additionally \(x-3x^2/2\le1/6\) in I, over the whole domain.

\paragraph{C182.8.2 Actual three-step mean and curvature costs}

In I, additional mean cost is at most \(37d/5\); the actual curvature coefficient is at most \(7d/4\), costing \(245d/72\). Hence loss \(607/40+19/4-37/5-245/72=821/90>9\).

In II, additional mean cost is at most \(134d/25\); coefficient \(13d/10\) costs \(169d/90\), giving \(64/5+23/4-134/25-169/90=10181/900>11\).

In III, additional mean cost is at most \(583d/120\); coefficient \(27d/20\) costs \(81d/40\), giving
\[
 35/4+27/5-583/120-81/40=109/15.                    \tag{C182.44}
\]
The negative-delta cost has been retained.

\paragraph{C182.8.3 Every finite length at least four}

In I and II, the tail norm is below \(6\sqrt d/25\), giving full vector norm below \(36\sqrt d/25\). Take \(h=3/8\); \(|\beta-hd|\le5d/8\), so curvature costs at most \(125d/288\). In I, \(2\delta_++2(D_1)_+\le3d/2\), giving loss
\[
 707/40-5508/625-125/288>42/5.                      \tag{C182.45}
\]
One may use \(5508/625<53/6\), \(125/288<7/16\), leaving \(2017/240>42/5\). In II, \(2\delta_++2(D_1)_+\le12d/25\), giving loss \(969/50-5508/625-125/288>10\).

In III retain its exact all-index gamma bound and take \(h=1/2\). The tail satisfies
\[
\begin{gathered}
\|B^{L-1}R\|\le[\theta^2+(2/9)\theta^{3/2}]\sqrt d\\
<(61/270)\sqrt d<(23/100)\sqrt d.
\end{gathered}
\]
The full vector norm is below \(143\sqrt d/100\), with norm cost \(347633d/40000<87d/10\). The positive scalar terms satisfy \(2\delta_++2(D_1)_+\le7d/10\), and curvature costs \(5d/18<7d/25\). The basic loss is
\[
 319/20-87/10-7/25=697/100>69/10.                   \tag{C182.46}
\]
This covers every finite connector, gamma index and terminal pair count.

\paragraph{C182.8.4 All five packets and every length range}


\begingroup\small
\setlength{\tabcolsep}{3pt}
\setlength{\LTleft}{0pt}\setlength{\LTright}{\fill}
\begin{longtable}{@{}p{0.2200\linewidth}p{0.7000\linewidth}@{}}
\caption{Five even multipair packets with the actual two-step root exception: parameter values.}\label{tab:parameters:q5-c182:1}\\
\toprule
Entry & \(L=1\) basic loss \\
\midrule\endfirsthead
\toprule
Entry & \(L=1\) basic loss \\
\midrule\endhead
\bottomrule\endfoot
\(a_1=2\) & \(185/24\); \(L=2\) basic loss: \(27/5\); \(L=3\) basic loss: \(32/5\); Every \(L\ge 4\) basic loss: \(49/5\). \\
\addlinespace[2pt]
\(a_1\ge 3\) & 12; \(L=2\) basic loss: 11; \(L=3\) basic loss: 12; Every \(L\ge 4\) basic loss: 14. \\
\addlinespace[2pt]
\(a_1=1,\allowbreak b_1\ge 3\) & 10; \(L=2\) basic loss: 10; \(L=3\) basic loss: 9; Every \(L\ge 4\) basic loss: \(42/5\). \\
\addlinespace[2pt]
\(a_1=1,\allowbreak b_1=2,\allowbreak a_2\ge 4\) & 12; \(L=2\) basic loss: 9; \(L=3\) basic loss: 11; Every \(L\ge 4\) basic loss: 10. \\
\addlinespace[2pt]
\(a_1=1,\allowbreak b_1=2,\allowbreak a_2=3\) & \(15/2\); \(L=2\) basic loss: \(13/2\); \(L=3\) basic loss: \(109/15\); Every \(L\ge 4\) basic loss: \(69/10\). \\
\end{longtable}
\endgroup


Each displayed value is a valid weaker paid rate, with the same additional \(64(j-2)/5\) in every cell. The five rows partition all original terminal lists, and the four columns cover every positive \(L\) in the same preselected direction. The minimum is \(27/5\), or \(13/2\) for the three minimum-one rows.
